\scanpage{29}
Donc il existe $u:T\to X$ tel que $u'=uf$. Alors on a \deleted{\unsure{que}} $\eta=F(u)(\xi)$, car en vertu de $\underset{\leftarrow}{L}_3$ a), il suffit de vérifier que
\[
F(f')(\eta)=F(f')F(u)(\xi)\quad\text{i.e.}\quad\eta'=F(u')(\xi),
\]
ce qui est vrai.
\hfill OK

\uline{N.B.} Il est plus naturel d'énoncer ces deux résultats ainsi: on a un hom de foncteurs $G\to F$ satisfaisant les conditions analogues pour $F$, et on \unsure{déduit|désire} que c'est un \uline{isom}.

\uline{Corollaire 2} Soient $S$ schéma artinien, $F:(\mathrm{Sch})_{/S}^{\circ}\mathbin{-}(\mathrm{Ens})$, pour que $F$ soit représentable par un $X$ loc. quasi-fini sur $S$, il f. et s. que $F$ satisfasse les conditions suivantes:

$1^\circ)$ $F$ est proreprésentable, les anneaux locaux qui interviennent dans la \unsure{proreprésentabilité} de $F$ sont \deleted{\unsure{finis}} \uline{artiniens}, i.e. finis sur \unsure{ceux} de $S$.

$2^\circ)$ $F$ satisfait les conditions $\underset{\leftarrow}{L}_1$, $\underset{\leftarrow}{L}_2$, $\underset{\leftarrow}{L}_3^{\unsure{00|\circ\circ}}$, $\underset{\rightarrow}{L}^{\unsure{00|\circ\circ}}\deleted{\unsure{4}}$
\mspage{28}
